\documentclass[11pt,a4paper]{article}

\usepackage[T1]{fontenc}
\usepackage[utf8]{inputenc}
\usepackage[turkish,english]{babel}
\usepackage{geometry}
\geometry{margin=2.5cm}
\usepackage{hyperref}
\hypersetup{colorlinks=true, linkcolor=blue, urlcolor=blue, citecolor=teal}
\usepackage{graphicx}
\usepackage{amsmath}
\usepackage{booktabs}
\usepackage{listings}
\usepackage{xcolor}
\usepackage{tikz}
\usetikzlibrary{positioning, arrows.meta, shapes.geometric}

\lstset{
  basicstyle=\ttfamily\footnotesize,
  frame=single,
  breaklines=true,
  columns=fullflexible,
  keywordstyle=\color{purple}\bfseries,
  commentstyle=\color{gray}\itshape,
  stringstyle=\color{teal},
}

\title{\textbf{OneProduct Agent OS}\\
\large Hermes Orchestration and OpenClaw Tool-Use\\
for End-to-End Single-Product E-Commerce}

\author{Gökhan Türkmen\\
\small \href{https://github.com/NeoMagic-coder/oneproduct-agent-os}{github.com/NeoMagic-coder/oneproduct-agent-os}}

\date{May 2026}

\begin{document}
\maketitle

\begin{abstract}
We describe \emph{OneProduct Agent OS}, a multi-agent platform that operates the
entire e-commerce lifecycle of a single product autonomously. The system is built
around two cooperating layers: \textbf{Hermes}, an orchestration plane that
decomposes user intent into a directed acyclic task graph; and \textbf{OpenClaw},
a permission-scoped tool-use plane that exposes 70+ self-describing tool
manifests. A heterogeneous roster of 18 specialised agents (research, brand,
pricing, growth, support, legal, etc.) cooperates through Hermes while sharing
the OpenClaw registry. We describe the lifecycle, confidence-based escalation,
and human-in-the-loop approval gates. A reference implementation ships as a
single-file Vite frontend (TypeScript) plus a FastAPI backend (Python) with a
mock LLM provider that gracefully upgrades to Google Gemini~\cite{gemini2024}.
\end{abstract}

\section{Introduction}

Modern e-commerce operations span at least twelve distinct surfaces---product
research, brand identity, store setup, listing optimisation, pricing,
advertising, organic content, fulfilment, customer support, reviews, analytics,
and legal compliance. Operating all of them for a single product typically
requires a multi-person team. Recent advances in foundation models and tool-use
\cite{schick2023toolformer, yao2023react} suggest that much of this surface can
be delegated to a cooperative set of specialised agents
\cite{wu2023autogen, shen2023hugginggpt}.

The contributions of this paper are:
\begin{enumerate}
  \item A two-layer architecture (\textbf{Hermes}~\textsection\ref{sec:hermes},
    \textbf{OpenClaw}~\textsection\ref{sec:openclaw}) that decouples
    orchestration from tool-use, mirroring the classical
    sense--plan--act separation but generalised to multi-agent settings.
  \item A confidence-aware lifecycle (\textsection\ref{sec:lifecycle}) with
    human-in-the-loop approval gates derived from risk classes.
  \item A pragmatic reference implementation released under MIT.
\end{enumerate}

\section{Architecture}
\label{sec:architecture}

\begin{figure}[h]
\centering
\begin{tikzpicture}[
  node distance=0.6cm and 1.0cm,
  agent/.style={draw, rounded corners, minimum width=2.6cm, minimum height=0.8cm, align=center, fill=blue!10},
  layer/.style={draw, dashed, rounded corners, inner sep=0.4cm},
  arrow/.style={-{Stealth[length=2mm]}, thick}
]
  \node[agent, fill=yellow!20] (user) {User};
  \node[agent, fill=orange!20, right=of user] (ceo) {CEO Agent};
  \node[agent, right=of ceo] (a1) {Research};
  \node[agent, below=of a1] (a2) {Pricing};
  \node[agent, right=of a1] (a3) {Brand};
  \node[agent, right=of a2] (a4) {Growth};
  \node[agent, fill=green!15, below=of user] (tools) {OpenClaw\\Registry};

  \draw[arrow] (user) -- (ceo);
  \draw[arrow] (ceo) -- (a1);
  \draw[arrow] (ceo) -- (a2);
  \draw[arrow] (a1) -- (a3);
  \draw[arrow] (a2) -- (a4);
  \draw[arrow] (a1) -- (tools);
  \draw[arrow] (a2) -- (tools);
  \draw[arrow] (a3) -- (tools);
  \draw[arrow] (a4) -- (tools);
\end{tikzpicture}
\caption{Hermes (top) decomposes a user request and routes it through a DAG of
specialised agents; each agent calls into the shared OpenClaw tool registry
(bottom).}
\label{fig:architecture}
\end{figure}

\subsection{Hermes orchestrator}
\label{sec:hermes}

Hermes is the single entry point for user-originated tasks. Given a message $m$,
Hermes performs three operations:
\begin{enumerate}
  \item \textbf{Intent routing}: a lightweight classifier produces a primary
    agent $a^\star$ and supporting agents $\mathcal{S}$. The reference
    implementation uses a keyword heuristic; production deployments substitute
    an LLM-based classifier.
  \item \textbf{Planning}: a directed acyclic task graph $G=(V,E)$ is
    constructed where each $v\in V$ is an (agent, payload) pair, and edges
    encode pre-requisites.
  \item \textbf{Execution \& merge}: ready nodes are executed in parallel; once
    all nodes resolve, an LLM provider produces a single executive summary
    constrained by a system prompt that injects the product context, routing
    rationale, and per-agent findings.
\end{enumerate}

\subsection{OpenClaw tool-use}
\label{sec:openclaw}

OpenClaw exposes tools as JSON Schema manifests with the following invariants:
\begin{itemize}
  \item Each manifest declares \texttt{allowed\_agents}; the executor rejects
    requests that violate this contract.
  \item Inputs and outputs are validated against the manifest schema before and
    after invocation.
  \item Calls run in a sandbox with a per-tool timeout and a maximum
    \texttt{retry.max\_attempts}. On exhaustion, a declared
    \texttt{fallback\_tool\_id} is attempted.
  \item Every execution emits a \texttt{ToolCallLog} for auditing and cost
    tracking; an optional per-context budget caps cumulative spend.
\end{itemize}

\begin{lstlisting}[language=Python, caption={OpenClaw permission check (excerpt).}]
def is_allowed(self, tool_id: str, agent_id: str) -> bool:
    tool = self.get(tool_id)
    if tool is None:
        return False
    if not tool.allowed_agents:
        return True
    return agent_id in tool.allowed_agents
\end{lstlisting}

\section{Lifecycle and approval}
\label{sec:lifecycle}

Tasks transition through the states
\texttt{created} $\rightarrow$ \texttt{triaged} $\rightarrow$ \texttt{assigned}
$\rightarrow$ \texttt{in\_progress} $\rightarrow$ \{\texttt{waiting\_tool\_result},
\texttt{waiting\_human\_approval}\} $\rightarrow$ \{\texttt{completed},
\texttt{failed}, \texttt{escalated}\}.

Each agent reports a confidence score $c \in [0,1]$. If $c$ drops below an
agent-specific threshold (default $0.6$), Hermes promotes the task to
\texttt{escalated} and notifies the operator. Actions that touch shared state
(pricing changes above $\pm 5\%$, mass mail send, legal text publication) are
flagged \texttt{requires\_approval} on the manifest and gate behind a
human-in-the-loop approval queue inspired by~\cite{ouyang2022rlhf}.

\section{Implementation notes}

The reference implementation is intentionally minimal:

\begin{itemize}
  \item \textbf{Frontend}: Vite 7 + React 19 + Tailwind v4, bundled to a single
    HTML file via \texttt{vite-plugin-singlefile} for trivial demos.
  \item \textbf{Backend}: FastAPI with Pydantic v2 schemas. Tool manifests live
    as JSON files under \texttt{apps/api/tools/manifests/}; the registry hot-loads
    them on boot.
  \item \textbf{LLM provider}: a thin abstraction over Gemini's
    \texttt{generateContent} endpoint with a deterministic mock fallback for
    offline development.
\end{itemize}

\section{Related work}

Toolformer~\cite{schick2023toolformer} introduced self-supervised tool use for a
single LLM; ReAct~\cite{yao2023react} extended this to interleaved
reasoning--action. AutoGen~\cite{wu2023autogen} and
HuggingGPT~\cite{shen2023hugginggpt} explored multi-agent collaboration but
without a separate tool plane. OneProduct Agent OS specifically targets the
e-commerce vertical and emphasises the operational concerns (approval, audit,
escalation) that are usually treated as afterthoughts.

\section{Limitations and future work}

The current router is keyword-based and brittle on novel phrasings. The mock
LLM provider is suitable for offline demos but does not reflect production
behaviour. We plan to (i) replace the router with a fine-tuned classifier; (ii)
add a streaming response path; and (iii) extend OpenClaw with live adapters for
Shopify, Trendyol, Amazon, and Meta Ads.

\bibliographystyle{plain}
\bibliography{refs}

\end{document}
